\chapter{The project deep-dive}
\label{ch:deep-dive}

\section*{What this chapter is for}

There is one round where the Staff decision is actually made, and it is not the
system design round. It is the conversation about something you have already
built \dash{} and its criteria are published, specific, and almost entirely
about scope rather than about technology.

\section{What it is scored on}

\reported{The criteria reported for this round are: whether you drove the
architectural decisions or executed somebody else's plan; whether you can
explain \emph{why}, over the alternatives; whether you can go several levels
deep when probed; and whether you assess your own work
honestly.}{field-guide-deep-dive}

And one more, about how the story is told: \reported{strong candidates frame
around impact \dash{} a number, a consequence \dash{} rather than around
technology names.}{field-guide-deep-dive}

\begin{kitnote}
Read that list against Chapter~\ref{ch:staff} and it is the same list. \emph{Did
you drive the decisions} is the unscoped-problem test. \emph{Why, over the
alternatives} is the judgement test. \emph{Several levels deep} is the
did-you-really-own-it test. \emph{Honest self-assessment} is the one nobody
prepares for and it is the one that separates people.

This is why the round decides the level. Everything else in the loop can be
passed by a very good Senior engineer.
\end{kitnote}

\section{Choosing the project}

Not the most impressive one. The one where you decided the most.

\begin{kittable}{@{}lXX@{}}
& \textbf{Weak choice} & \textbf{Strong choice} \\
\midrule
Scope & you built a well-specified component & you decided what the problem was \\
Alternatives & one obvious approach & a real fork with a real cost either way \\
Depth & you know your layer & you know why the layer below behaved as it did \\
Outcome & it shipped & something measurable changed, or it did not and you know why \\
Age & last month, still unresolved & finished, with consequences you have seen \\
\end{kittable}

\judgement{A moderately sized project you owned end to end beats a large one
where you owned a component, every time, and candidates reliably choose the
large one because it sounds better. The large one produces an answer to ``who
decided that?'' which is somebody else's name, and that answer is the end of
the Staff conversation.}

The project does not have to be from this domain. Scope, ambiguity and
influence do not belong to a technology, and an unscoped problem you solved in
a previous life is better evidence than a well-specified one you solved last
week.

\section{The structure that survives probing}

\begin{kitmethod}
\textbf{The situation, in two sentences.} What was wrong, and for whom. Not the
architecture.

\textbf{Why it was your problem.} How it came to you, and how much of the
scoping was yours. If somebody handed you a specification, say so \dash{} and
then say what you decided inside it, because there is always something.

\textbf{The fork.} The one decision that mattered, both options, the cost of
each, and \emph{what you knew at the time}. The last part is what makes it a
decision rather than a retelling.

\textbf{What would have changed your mind.} Prepare this sentence. It is the
strongest single thing you can say in this round, and it is the one that
distinguishes somebody who decided from somebody who preferred.

\textbf{What happened.} With a number if you have one.

\textbf{What you got wrong.} Genuinely wrong, with what you would do
differently. This is the honest-self-assessment criterion and it is being
scored.
\end{kitmethod}

\section{Going several levels deep}

You will be probed downwards until you stop being able to answer. That is the
instrument, and it is fine to reach the bottom \dash{} what is being measured is
where it is and how you behave when you get there.

So: know the layer below the one you worked in. If you built on a database,
know why the query planner did what it did. If you built on a queue, know its
delivery guarantee and what you relied on. If you built on a model provider,
know what you assumed about its behaviour under load.

And when you reach the bottom, say so plainly. \judgement{``I do not know, and
here is how I would find out'' is a strong answer at Staff and a weak one at
Junior, which is worth knowing, because the instinct built up through earlier
interviews is to keep talking.}

\section{The honesty criterion}

The published criteria name honest self-assessment explicitly, and it is the
least-prepared item on the list.

What it looks like in practice: naming the part of the design you are least
confident about, without being asked; distinguishing what you measured from
what you believe; saying which of the outcomes you can attribute to your
decision and which you cannot; and not claiming an improvement that had three
causes as though it had one.

\begin{kittrap}
The failure mode is not dishonesty. It is the accumulated smoothing that
happens when a story has been told a dozen times: the alternatives get
simplified, the uncertainty falls away, and the outcome attaches itself to your
decision. Nobody decides to do this.

The remedy is to go back to the source \dash{} the design document, the
tickets, the messages from the time \dash{} before you tell the story to an
interviewer. You will usually find that the decision was less clean and more
interesting than the version you have been telling.
\end{kittrap}

\begin{drill}
Fill in \code{templates/project-deep-dive.md} for your strongest project,
against the six-part structure. Then have somebody ask you ``why?'' five times
in a row, each time about your previous answer. Note where you run out. That
point is where you need to go and read something, and it is almost always one
layer below where you expected.
\end{drill}

\section*{What this chapter does not claim}

The criteria are reported from one corpus of disclosed processes. They are
consistent with what published ladders say the level means, which is corroboration
of a sort and not independent evidence.

The claim that this round decides the level more than the system design round
does is the author's judgement. The argument for it is in the note above: the
other rounds can be passed by a very good Senior engineer, and this one is the
one that asks the question the ladder actually turns on.
